\clearpage
\section{Index of principal notation}\label{sec:symbolic-index}

The entries refer to the conventions at the indicated places. Symbols used
only in a local calculation retain the meaning declared there. In particular,
a change of type, scale, normalization, or limiting order requires the stated
bridge; repeated letters alone do not identify mathematical objects.

\begin{description}
\item[$G$, $r_G$, $\mathfrak g$] The fixed compact matrix group in its
specified representation of complex dimension $r_G$, and its real Lie
algebra (Section~\ref{ex-statement}).
\item[$x=(\boldsymbol x,t)$, $\overline V_+$] Euclidean coordinates with
time last, and the closed forward cone for the Minkowski form
$t^2-|\boldsymbol x|^2$ (Section~\ref{ex-statement}).
\item[$P_p$, $D_N$, $a_N$] Physical period, lattice denominator, and mesh,
with $a_N=D_N^{-1}$. The site count per axis is $P_pD_N$
(Section~\ref{ex-statement}).
\item[$\Lambda_{N,p}$, $U$, $U_{ij}(x)$] Periodic sites, a group-valued
link configuration, and its oriented plaquette holonomy
(Section~\ref{ex-lattice}).
\item[$S_\ell$, $\beta$, $\mu_{\ell,\beta}$] The finite plaquette action,
its bare coefficient, and the normalized Wilson probability. Each pair
of axes is counted once and the trace is divided by $r_G$
(Section~\ref{ex-lattice}). This action is distinct from the correlation
table $S$ below.
\item[$O_{N,p,v}$, $M_{N,p}$, $W_{N,p}$] A token insertion, its finite
moment, and the smeared moment. The latter includes $a_N^{4n}$ at arity
$n$ \eqref{ex-smearing}.
\item[$\operatorname{Conf}_n$, $\mathrm{Gen}$, $\mathrm{Cyc}_p$,
$\mathrm{NoHit}_p$] Distinct-point configurations, their generic subset
with separated coordinate projections, and the two period-dependent
configuration domains defined in Section~\ref{sec:ex-limits}; the
coordinatewise conditions are not interchangeable.
\item[$f_{\rm mesh}$, $f_{\rm per}$] The two increasing sequences of an
internal extraction order. The ordered tail has a period threshold outside
a mesh threshold that may depend on the period
(Definition~\ref{def:ex-limits-estimates}).
\item[$S$, $S^\chi$, $S^{\rm mix}$] The limiting Euclidean table, a
character component, and its mixed correlation. The mixture is the integral
of the character correlations in \eqref{ex-mixture}; it is the mixed
correlation that appears in the Wilson limit.
\item[$\mathcal F_d$, $A$] In the public target, the finite-dimensional
space at grade $d$ and a field label. A grade is not automatically a
scaling dimension (Section~\ref{ex-statement}).
\item[$X$, $\nu$, $\chi_*$] In the public target, a compact phase space,
its full-support probability, and a selected character. The symbols are
local to that discussion; the construction uses $\Xi$ and $\mu$ in
Proposition~\ref{ex-character-measure}.
\item[$\mathcal H$, $\Omega$, $\mathcal D$, $\Phi$] The selected theory's
Hilbert space, vacuum, common invariant domain, and fields
(Definition~\ref{ex-reconstructed}). The domain is distinct from any
locally named contact-device or matrix object.
\item[$A$, $M$, $P_{\mathcal H}$ in the sector calculation] The completed
precision $A=M^*M+P_{\mathcal H^c}$, selected row map, and projection
onto active bonds \eqref{eq:sector-precision}. The Gaussian covariance,
when defined, is the inverse precision; the completion is identity on
inactive coordinates.
\item[$D_{v,a}$ and $\mathscr D_{v,a}$ in the free-kernel calculation]
The train drive and level response, respectively. They have different
formulas and cannot be substituted for each other
(Section~\ref{sec:ex-lattice-free-kernel-components}).
\item[$A$, $A_0$, $C$ in Gaussian reassembly] The completed sector
precision, flat reference precision, and localized positive finite-range
covariance. The correction $C^{-1}-A$ is retained; $C$ is not simply
identified with $A^{-1}$ (Section~\ref{sec:gaussian-reassembly}).
\item[$a_m,b_m,c_m$ in background paths] The quadratic-loop, cubic,
and quartic coefficient arrays at a fixed scale, in the prescribed
real coefficient spaces (Section~\ref{sec:background-path-convergence}).
This $a_m$ is local and differs from the physical mesh.
\item[$\varepsilon$ in the healthy-field comparison] The field error and
per-step row/covariance error are separate quantities in the inputs to the
comparison. The stated conclusion uses a common bound for both
(Section~\ref{sec:marked-healthy-cumulant-launch}).
\item[$\lambda$ in the sector calculation] The normalization $q^{-1/2}$
for $q$ signed channels, so its square gives the orientation average
(Section~\ref{sec:sector-gram}). Other local uses of $\lambda$ are
introduced separately.
\item[$\lambda_G$, $b_0(G)$, $s_G$] The annulus coefficient,
canonical running coefficient, and shooting-line slope, respectively.
The two Casimir definitions use the same color inner product, giving
$\lambda_G=2b_0(G)$. The match to $s_G$ is a separate input (Section~\ref{sec:surviving-interval-shooting}).
\item[$E_d$, $\mathcal F_d$, $\pi_d$, $\tau_{I,D}$ in the annulus]
Reference label space, selected-table field space, quotient map, and
spectator tester (Section~\ref{sec:ex-lattice-actual-stress-annulus}).
The finite-volume labels and their quotient fields have distinct domains.
\item[$\mathsf P_r(M;C)$ in the annular estimate] A bound on the matrix path and its first two
derivatives with weights $r_i^j$ for $j=0,1,2$
(Section~\ref{sec:ex-annulus-fleet-faces}). Here $r_i$ is a local price
scale, not automatically an annulus or short-distance radius.
\item[$T(r)$] The selected two-point trace quantity in the final
short-distance requirement \eqref{ex-trace-limit}; this is a separate
condition from the operator-product estimate in \eqref{ex-ope-remainder}.
\end{description}
